% UML 2.5 use case diagram (1 of 2) — learning and practice. Build of 2026-11-25.
% Compile: lualatex usecase_learning.tex (standalone PDF, included in the report).
% Grid (cm): actors x=0; column A x=4.2 (use cases an actor starts, in the order they
% happen); column B x=10.0 and C x=13.9 (use cases that extend another, placed below it).
% Rows 1.5 cm apart. All use cases are drawn alike; their implementation status is in the
% report's status table.
\documentclass[tikz, border=3mm]{standalone}
% Shared TikZ styles for the UML 2.5 use case diagrams of the report.
% Conventions followed (Ambler, The Elements of UML 2.0 Style; Microsoft UML use case
% guidelines): actors are roles drawn outside the system boundary, primary actors top-left,
% a specialised actor below its parent with the generalization arrow pointing to the parent
% (shared-target style: one hollow triangle under the parent, one trunk, one branch per
% child); use cases are horizontal ellipses named with a strong verb, stacked top to bottom
% in the order they happen; an included use case sits to the right of its base, an extending
% use case below its base; associations carry no arrowheads. All use cases are drawn alike;
% the implementation status is given in the report's status table, not in the diagram.
% Layout is placed by hand on a grid.
\usetikzlibrary{arrows.meta,shapes.geometric,positioning,calc,fit,backgrounds}

\definecolor{umlline}{HTML}{1F2937}

\tikzset{
  every node/.style={font=\sffamily\footnotesize},
  % a use case: an ellipse that hugs its (at most two-line) name
  uc/.style={ellipse, draw=umlline, line width=0.5pt, fill=white, align=center,
             inner xsep=0.3mm, inner ysep=0.2mm, minimum height=8mm, minimum width=26mm},
  assoc/.style={draw=umlline, line width=0.5pt},
  gen/.style={draw=umlline, line width=0.5pt, -{Triangle[open, length=2.6mm, width=2.6mm]}},
  dep/.style={draw=umlline, line width=0.5pt, dashed, -{Straight Barb[length=2mm, width=2.2mm]}},
  deplabel/.style={font=\sffamily\scriptsize, fill=white, inner sep=0.8pt},
  boundary/.style={draw=umlline, line width=0.7pt},
}

% \actor{node name}{position}{label}: a stick figure; the node is the figure's body.
\newcommand{\actor}[3]{%
  \begin{scope}[shift={#2}, line width=0.6pt, draw=umlline]
    \draw (0,0.34) circle (0.13);
    \draw (0,0.21) -- (0,-0.14);
    \draw (-0.23,0.08) -- (0.23,0.08);
    \draw (0,-0.14) -- (-0.18,-0.42);
    \draw (0,-0.14) -- (0.18,-0.42);
    \node[inner sep=0pt, minimum width=5mm, minimum height=10mm] (#1) at (0,0) {};
    \node[below=0.46cm of #1.center, align=center, font=\sffamily\footnotesize] (#1-label) {#3};
  \end{scope}%
}

% \specialize{child}{parent}{trunk x}: generalization from a child actor (placed to the
% right of the trunk) to its parent, drawn orthogonally: a horizontal branch from the child
% to the trunk, then up the trunk to the hollow triangle under the parent's name. Branches
% of several children share the trunk and the triangle (UML shared-target style).
\newcommand{\specialize}[3]{%
  \draw[gen] ($(#1.west)+(0.02,0)$) -- (#3,0 |- #1.west) -- ($(#2-label.south)+(0,-0.05)$);%
}


\begin{document}
\begin{tikzpicture}

% ---------------- base use cases: column A, in time order ----------------
\node[uc] (UC00) at (4.2,  0.00) {UC-00 Register\\an Account};
\node[uc] (UC01) at (4.2, -1.50) {UC-01 Log In};
\node[uc] (UC02) at (4.2, -3.00) {UC-02 Analyze\\a Song};
\node[uc] (UC07) at (4.2, -4.50) {UC-07 Manage the\\Library and Playlists};
\node[uc] (UC04) at (4.2, -6.00) {UC-04 Practice\\a Song};
\node[uc] (UC19) at (4.2, -8.10) {UC-19 Report a Chord\\Error and Send Feedback};
\node[uc] (UC22) at (4.2, -9.60) {UC-22 Export\\a Chord Chart};
\node[uc] (UC03) at (4.2,-11.10) {UC-03 Detect Chords\\from a MIDI Instrument};
\node[uc] (UC06) at (4.2,-12.60) {UC-06 Practice\\with the Microphone};

% ---------------- extending use cases: below the use case they extend ----------------
\node[uc] (UC05) at (10.0,-6.00) {UC-05 Edit\\a Chord};
\node[uc] (UC18) at (10.0,-7.00) {UC-18 Use\\the Metronome};
\node[uc] (UC20) at (10.0,-8.00) {UC-20 Correct the Rhythm\\and Song Sections};

% ---------------- system boundary ----------------
\begin{scope}[on background layer]
  \node[boundary, fit=(UC00)(UC06)(UC20), inner xsep=4mm, inner ysep=5mm,
        label={[anchor=north, font=\sffamily\bfseries\footnotesize]north:ChordSense Pro}] (SYS) {};
\end{scope}

% ---------------- actors (primary actors top-left, specialised actor below its parent) ----
\actor{Guest}{(0, 0.00)}{Guest}
\actor{Reg}{(0,-1.50)}{Registered User}
\actor{Free}{(0,-5.10)}{Free User}
\actor{Paid}{(0,-11.10)}{Paid User (Pro)}

\draw[gen] ($(Free.north)+(0,0.05)$) -- ($(Reg-label.south)+(0,-0.05)$);
\draw[gen] ($(Paid.north)+(0,0.05)$) -- ($(Free-label.south)+(0,-0.05)$);

% ---------------- associations (no arrowheads) ----------------
\draw[assoc] (Guest.east) -- (UC00.west);
\draw[assoc] (Reg.east)   -- (UC01.west);
\draw[assoc] (Free.east)  -- (UC02.west);
\draw[assoc] (Free.east)  -- (UC07.west);
\draw[assoc] (Free.east)  -- (UC04.west);
\draw[assoc] (Free.east)  -- (UC19.west);
\draw[assoc] (Paid.east)  -- (UC22.west);
\draw[assoc] (Paid.east)  -- (UC03.west);
\draw[assoc] (Paid.east)  -- (UC06.west);

% ---------------- «extend»: from the extending use case to the base use case ----------------
\draw[dep] (UC05.west) -- node[deplabel, pos=0.5] {\guillemotleft extend\guillemotright} (UC04.east);
\draw[dep] (UC18.west) -- node[deplabel, pos=0.45] {\guillemotleft extend\guillemotright} (UC04.-14);
\draw[dep] (UC20.west) -- node[deplabel, pos=0.35] {\guillemotleft extend\guillemotright} (UC04.-30);

\end{tikzpicture}
\end{document}
